% ============================================================
% UPDATED SECTION 2.4 + PROPOSITION 1 (β₁ stripped)
% This file is intended as a drop-in replacement for earlier drafts.
% ============================================================

\subsection{Model space and revision graph}\label{sec:model_space}

The preceding sections characterize DLN stages by the structure of the agent's
belief-dependency graph~$G$.  This section introduces a second formal object:
the \emph{model space}~$\mathcal{M}$ of candidate belief structures and the
\emph{revision graph} $\mathcal{R}=(\mathcal{M},\mathcal{T})$ whose edges are
transitions the structural learning cycle can execute.  Together, $G$ and
$\mathcal{R}$ provide a two-object characterization: $G$ determines the cost of
inference within a representation (Level~1), while $\mathcal{R}$ determines the
agent's capacity to evaluate and revise that representation (Level~2).

\paragraph{Definition (model space).}
Let $\mathcal{M}$ be the set of belief-dependency structures available to the
agent.  In our instantiation,
\[
\mathcal{M} = \bigl\{G_{\mathrm{tab}},\; G_{F}\bigr\},
\]
where $G_{\mathrm{tab}}$ is the tabular structure (null graph on $K$ nodes,
maintaining independent per-option beliefs) and $G_F$ is the factor structure
(bipartite DAG with $F$ shared-factor nodes connecting to $K$ option nodes).
Each $M \in \mathcal{M}$ determines (i)~the agent's representational state size
($\Theta(K)$ for $G_{\mathrm{tab}}$, $\Theta(F)$ for~$G_F$) and (ii)~the
inference procedure used for action selection.

\paragraph{Definition (revision graph).}
The \emph{revision graph} $\mathcal{R} = (\mathcal{M}, \mathcal{T})$ is a
directed graph whose vertices are elements of $\mathcal{M}$ and whose directed
edges $(M, M') \in \mathcal{T}$ represent transitions the structural learning
cycle can execute.
In the expand-only implementation (Section~\ref{sec:slc}), the transition set is
\[
\mathcal{T}_{\text{expand}} = \bigl\{(G_F,\, G_{\mathrm{tab}})\bigr\}.
\]
The agent can expand from the factor model to the tabular model when the
predictive test detects inadequacy, but it cannot contract back.  This revision
graph is a single directed edge: $G_F \to G_{\mathrm{tab}}$.

\paragraph{Completing the cycle: contraction.}
We augment the structural learning cycle with a \emph{contraction} mechanism
(Section~\ref{sec:contraction}).  While operating under $G_{\mathrm{tab}}$, the
agent periodically evaluates whether a compressed factor hypothesis is again
adequate.  If the predictive test favors the compressed model, the agent
executes a return transition $(G_{\mathrm{tab}}, G_F)\in\mathcal{T}$ and resumes
$O(F)$-scaled computation.

\paragraph{Revision capacity by DLN stage.}
The revision graph characterizes each stage's meta-cognitive capacity in
graph-combinatorial terms:
\begin{itemize}[leftmargin=*]
\item \textbf{Dot.}  No model space ($|\mathcal{M}|=0$); $\mathcal{R}$ is empty.
  No revision capacity.
\item \textbf{Linear.}  $\mathcal{M} = \{G_{\mathrm{tab}}\}$; a single vertex
  with no edges.  The agent occupies a fixed point in model space.
\item \textbf{Network (expand-only).}  $\mathcal{M} = \{G_F, G_{\mathrm{tab}}\}$;
  one directed edge $G_F \to G_{\mathrm{tab}}$.  The agent can detect model
  failure and expand, but cannot return to a compressed representation.
\item \textbf{Network (full cycle).}  Two directed edges (expand and contract).
  $\mathcal{R}$ contains a cycle: the agent can depart from a model, operate
  under an alternative, and return.  This is the graph-theoretic signature of
  full revision capacity.
\end{itemize}
The functional consequence of the return transition is stated in
Proposition~\ref{prop:crossover}(iii): an agent that can return to $G_F$ recovers
$O(F)$ cost scaling in bounded time; an agent that cannot is permanently locked
at $O(K)$.

\paragraph{Two levels of computation.}
Level~1 computation concerns inference \emph{within} a chosen belief structure
$G$ (e.g., updating and using $\widehat{m}_f$ and $\widehat{b}_f$ under $G_F$ or
updating and using $\widehat{\mu}_i$ under $G_{\mathrm{tab}}$).
Level~2 computation concerns monitoring model adequacy and executing transitions
in $\mathcal{R}$ (hypothesis testing, expansion, and contraction).
Proposition~\ref{prop:crossover}(i) shows that Level~2 is only cost-effective
when the option set is sufficiently large relative to the factor dimension, and
Proposition~\ref{prop:crossover}(iii) shows that including a return transition
enables bounded recovery after model failure.

\paragraph{Relationship to tractable-inference results.}
Classical work on probabilistic graphical models \cite{Pearl1988} and the
treewidth / tractable-cognition literature \cite{vanRooij2008,Kwisthout2010}
analyze the complexity of inference \emph{within a fixed model structure} and
establish that low-treewidth graphs enable efficient inference.  The within-step
belief graph for our Network agent is itself a low-treewidth bipartite DAG; we
do not claim that cyclic belief graphs make inference cheaper than trees.
Instead, we analyze a representation-construction question: when the world
admits shared structure with $F\ll K$ and the factor hypothesis is empirically
adequate, representing shared structure reduces the number of distinct
quantities that must be learned from $K$ to $F$; when structure is absent or
incorrect, expansion toward a tabular representation removes the advantage.

\begin{proposition}[Cost separation and recovery bound]\label{prop:crossover}
Let an agent face a $K$-option decision environment with $F$-factor structure
($F \leq K$).  Let $c_{\mathrm{param}}$ be the per-element parameter learning
cost and $c_{\mathrm{meta}}$ be the Level~2 model-revision overhead (navigation
of the revision graph $\mathcal{R}$, Section~\ref{sec:model_space}).  Then:
\begin{enumerate}[label=(\roman*)]
\item \textbf{Crossover.}
For $K < K^* = F + c_{\mathrm{meta}}/c_{\mathrm{param}}$, Linear achieves lower
total cognitive cost than Network (Network cannot do better than operating in a
correct compressed model and still paying its Level~2 overhead).
For $K > K^*$, Network under a correct factor hypothesis achieves lower total
cost by replacing $O(K)$ per-episode parameter maintenance with $O(F)$, up to the
fixed overhead $c_{\mathrm{meta}}$.

\item \textbf{Catastrophic failure under cross-terms.}
When the environment contains stakes-factor structure with covariance coupling,
Linear's utility degrades without bound in cumulative exposure $E_t$ because it
cannot represent the marginal penalty term $2E_t b_i$.  Network's utility
remains bounded because it tracks $E_t$ and computes marginal impact.

\item \textbf{Recovery bound.}
When the factor hypothesis is incorrect and the agent expands to
$G_{\mathrm{tab}}$, an agent whose revision graph $\mathcal{R}$ includes a
return transition $(G_{\mathrm{tab}}, G_F) \in \mathcal{T}$ recovers $O(F)$ cost
scaling within $n_{\mathrm{contract}} + w + n_{\mathrm{converge}}$ steps after
factor structure becomes present, where $n_{\mathrm{converge}}$ is the number of
observations required for within-factor variance to fall below the contraction
threshold.  An agent without a return transition (expand-only Network or Linear)
remains at $O(K)$ cost indefinitely.
\end{enumerate}

Parts (i) and (ii) are tested by the existing simulation results.
Part (iii) is tested by the contraction ablation
(Section~\ref{sec:contraction}).
\end{proposition}

\paragraph{Proof sketch for (i).}
Linear maintains $K$ independent estimates, so its per-episode parameter cost is
$c_{\mathrm{param}}K$.  Network operating under a correct compressed hypothesis
maintains $F$ shared estimates, giving $c_{\mathrm{param}}F$, but additionally
pays a fixed revision overhead $c_{\mathrm{meta}}$ for hypothesis monitoring and
possible model transitions.  Comparing $c_{\mathrm{param}}K$ to
$c_{\mathrm{param}}F+c_{\mathrm{meta}}$ yields the crossover
$K^* = F + c_{\mathrm{meta}}/c_{\mathrm{param}}$.

\paragraph{Proof sketch for (ii).}
With quadratic exposure penalty $\Delta P_t(i)=\lambda(2E_t b_i + b_i^2)$, the
marginal stakes impact depends on the sign interaction between $E_t$ and $b_i$.
A policy that ignores the cross-term cannot select actions that reliably reduce
$|E_t|$, so $E_t$ behaves like an uncontrolled random walk and the cumulative
quadratic penalty grows without bound.  Tracking $E_t$ and using the true
marginal form allows Network to select hedging actions that keep $E_t$ near
zero, bounding the penalty.

\paragraph{Proof sketch for (iii).}
After expansion to $G_{\mathrm{tab}}$, the agent periodically evaluates whether a
compressed factor hypothesis is again adequate (Section~\ref{sec:contraction}).
If $\mathcal{R}$ contains a return transition, then once factor structure is
present the predictive test will favor the compressed model within an evaluation
window of length $w$, after at most $n_{\mathrm{contract}}$ steps until the next
test begins, plus $n_{\mathrm{converge}}$ observations needed for within-factor
variance to drop below the contraction threshold.  Without a return transition,
the agent cannot exit $G_{\mathrm{tab}}$ and remains locked at $O(K)$ scaling.
